% Better Poster (Morrison-style) — A0 landscape, big-finding layout.
\documentclass[a0paper,fleqn]{betterposter}

\begin{document}
\betterposter{
	%%%%%%%% MAIN COLUMN

	\textbf{The proposed method cuts error rate by 30\%}
	\\over the strongest baseline,
	\\while using \textbf{half the training data}.
}{
	%%%%%%%% LEFT COLUMN

	\title{A Legible Result Communicated Over Coffee}
	\author{A. Researcher}
	\author{B. Collaborator}
	\institution{Department of Computational Something, University of Somewhere}

	\section{Introduction}
	This poster is built on Rafael Bailo's Better Poster template, which
	implements Mike Morrison's \#betterposter layout: the main finding
	stated in giant text in the centre, with supporting detail relegated
	to narrow side columns. Replace this section with the motivation for
	the work.

	\section{Methods}
	\begin{itemize}
		\item Describe the approach, dataset, or apparatus here.
		\item A second bullet, for a second step of the method.
		\item A third, for whatever completes the picture.
	\end{itemize}

	\section{Result Table}
	\begin{tabular}{lcc}
		Condition & Metric A & Metric B \\
		\hline
		Baseline  & 0.42     & 12.3     \\
		Proposed  & 0.61     & 9.7
	\end{tabular}

	\section{Conclusion}
	One or two sentences stating the takeaway a passer-by should leave
	with, even if they read nothing else on the poster.
}{
	%%%%%%%% RIGHT COLUMN

	\section{Discussion}
	Summarise what the main finding means, and how it compares to prior
	work. This is also the natural place for limitations or threats to
	validity.

	\section{References}
	[1] A. Author, ``A Foundational Paper,'' Journal of Examples, 2020.

	[2] B. Author and C. Author, ``A Follow-up Study,'' Proc. of the
	Conference on Placeholders, 2022.
}
\end{document}
